\documentclass[12pt]{article}

\usepackage[a4paper,margin=1in]{geometry}
\usepackage[T1]{fontenc}
\usepackage[utf8]{inputenc}
\usepackage{lmodern}
\usepackage{setspace}
\usepackage{parskip}
\usepackage{enumitem}

\setstretch{1.15}

\begin{document}

\begin{center}
    {\LARGE \textbf{Economics Is Everywhere}}\\[0.5em]
  
\end{center}

Good evening everyone,

Thank you so much for being here. It is a real pleasure to speak with both students and parents tonight.

[Pause]

Let me start with a simple question.

When you hear the word \textbf{economics}, what is the first thing that comes to mind?

[Pause for audience reaction]

Money?\\
Inflation?\\
Banking?\\
Graphs?\\
Maybe... stress?

[Smile]

For some people, economics sounds exciting.\\
For others, it sounds like something that comes with too many charts and not enough happiness.

But tonight, I want to show you that economics is actually much more interesting, much more practical, and much more connected to everyday life than most people think.

[Pause]

In fact, economics is not just about money.

It is about \textbf{choices}.\\
It is about \textbf{decision-making}.\\
It is about \textbf{how people, families, businesses, and governments deal with limited resources}.

And that means economics is not just in textbooks.\\
It is in your kitchen.\\
It is in your shopping cart.\\
It is in your bank account.\\
It is in your career choices.\\
It is in housing prices.\\
It is in social media ads.\\
And yes, it is even in that moment when you say,\\
“I should save money this month”\\
...and then order takeout three times in one week.

[Pause for laughter]

Tonight, I want to do a few things.

I want to explain what economics is in plain language.\\
I want to show you how it affects everyday life.\\
I want to introduce some really interesting areas like \textbf{behavioural economics} and \textbf{neuroeconomics}.\\
And I want to end by talking about why economics is worth studying and why it can lead to many exciting career opportunities.

So let us begin with the most basic question.

\section*{What is economics?}

At its core, economics is the study of how people make choices when they cannot have everything they want.

[Pause]

That is really it.

In economics, we call this \textbf{scarcity}.

Scarcity does not necessarily mean extreme poverty or shortage. It simply means that resources are limited.

Your time is limited.\\
Your money is limited.\\
Your energy is limited.\\
Even your attention is limited.

And because resources are limited, choices have to be made.

Let me ask the students here:

How many of you would like to have more sleep, better grades, more free time, more money, and less stress?

[Pause for hands / reaction]

Exactly.

The problem is that there are only 24 hours in a day.

That is economics.

Let me ask the parents:

How many of you would like lower grocery bills, lower gas prices, affordable housing, more savings, and maybe a vacation too?

[Pause]

Yes. Same problem.

Resources are limited. Choices must be made.

That is economics too.

So economics starts with a very human reality:\\
\textbf{we cannot have everything at once.}

And because of that, every choice has a cost.

\section*{Opportunity cost: the hidden cost of every decision}

One of the most important ideas in economics is called \textbf{opportunity cost}.

That sounds technical, but it is actually very simple.

Opportunity cost is the value of the next best thing you give up when you make a choice.

So if a student spends three hours watching videos instead of studying, the cost is not just the three hours.

The cost is also what those three hours could have produced:
better understanding, a better grade, less panic later.

[Pause]

And yes, economics has a way of making procrastination sound very academic.

If you stay up late scrolling through your phone instead of sleeping, economists call that a trade-off.

Your parents call it something else.

[Pause for laughter]

This idea applies everywhere.

If a family spends more money on one thing, it has less for something else.\\
If a government spends more in one area, it may have to cut back in another.\\
If you choose one university program, you are also choosing not to pursue another path.

Economics trains us to ask:\\
What am I gaining?\\
What am I giving up?\\
And is it worth it?

That is a very useful life skill.

\section*{Economics in everyday life}

Sometimes people imagine economics is something far away, happening only in government offices or stock markets.

But economics is everywhere.

Let me give you a few examples.

\subsection*{Coffee shop example}

Imagine the campus coffee shop at 8:30 in the morning.

There is a long line.

Why?

Because many students want coffee at the same time, and the coffee shop can only serve so many people at once.

That is scarcity.\\
That is supply and demand.\\
That is economics.

Now students may call it “a terrible morning.”

Economists call it “an allocation problem.”

[Pause for laughter]

\subsection*{Concert tickets}

Think about concert tickets.

A famous singer announces a concert.\\
Tickets sell out in minutes.

Why?

Because demand is high, but the number of seats is limited.

Then what happens?

Resale prices go up.\\
A \$100 ticket becomes \$300.

Now we are talking about scarcity, willingness to pay, market behaviour, and incentives.

Economics again.

\subsection*{Grocery shopping}

Now let us move to something every family understands: grocery shopping.

You go into the store for five things.\\
You come out wondering whether strawberries now count as luxury goods.

[Pause for laughter]

That is where inflation becomes real.

Inflation means prices are rising over time, and your money buys less than before.

You may hear about inflation on the news.\\
But you truly understand it when the grocery bill somehow becomes shocking every single week.

Economics helps explain why that happens and what it means for families, workers, and businesses.

\subsection*{Housing}

Now think about housing.

Why is housing so expensive in many places?

That is one of the biggest questions people ask today.

Economics helps us think about the possible reasons:
population growth, limited supply, construction costs, borrowing costs, regulations, and demand pressures.

Economics does not always give us one easy answer.\\
But it does help us ask the right questions.

And that matters.

\section*{Economics helps us understand big issues}

Economics is especially useful because it helps us make sense of major issues in society.

For example:

Why does inflation rise?\\
Why do interest rates go up?\\
Why do some people struggle to find jobs?\\
Why do wages differ so much across occupations?\\
Why are some countries richer than others?\\
Why do governments tax and spend the way they do?

These are not just academic questions.\\
These are real-life questions.

And economics gives us a way to think through them carefully.

One of the most useful habits economics teaches is this:

Do not just ask, “What happens immediately?”\\
Also ask, “What happens next?”

That is very important.

Because sometimes a policy sounds great at first, but creates unexpected consequences later.

Economics teaches us to look beyond the first effect and think about the second and third effects too.

That is a very powerful habit --- not just in school, but in life.

\section*{Economics is about people}

Now I want to emphasize something important.

Economics uses data.\\
It uses graphs.\\
It uses models.

But economics is really about people.

It is about families trying to manage a budget.\\
It is about students choosing a future.\\
It is about workers looking for stable jobs.\\
It is about businesses making decisions.\\
It is about governments trying to improve people’s lives.

So economics is not cold.\\
Good economics is deeply human.

It tries to understand how choices affect well-being, opportunity, fairness, and living standards.

\section*{Interesting fields in economics}

Now let us get to one of the most fun parts: the different branches of economics.

A lot of people think economics is just one narrow subject.

It is not.

There are many fascinating fields, and some of them are much more exciting than people expect.

\subsection*{Behavioural economics}

Let me start with one of the most popular ones: \textbf{behavioural economics}.

Traditional economics often assumes that people make rational decisions.

Behavioural economics says:

“Well... sometimes. But not always.”

[Pause]

Behavioural economics studies how people actually behave.

Why do people procrastinate?\\
Why do we overspend?\\
Why do we make emotional decisions?\\
Why do we keep saying we will start tomorrow?

This field recognizes that human beings are not robots.

We are influenced by habits, emotions, peer pressure, fear, overconfidence, and temptation.

For example, many people know they should save more money.

They also know they should sleep earlier, eat better, and spend less time on their phones.

And yet...

[Pause]

here we are.

[Pause for laughter]

Behavioural economics studies that gap between what people intend to do and what they actually do.

This field is used in real life to improve retirement saving, health decisions, public programs, and consumer protection.

\subsection*{Neuroeconomics}

Now here is one that sounds very futuristic: \textbf{neuroeconomics}.

Neuroeconomics is where economics meets psychology and neuroscience.

It studies what happens in the brain when people make decisions.

For example:
How do people respond to risk?\\
Why do rewards feel rewarding?\\
How do emotions affect our choices?\\
Why are some people more impulsive than others?

So yes, economics has expanded into studying the brain.

Now, that does not mean economists can read minds.

Let us be realistic.

Economists still cannot predict what a family wants for dinner when everyone says, “Anything is fine.”

[Pause for laughter]

But neuroeconomics helps us understand decision-making at a deeper level, especially in areas like trust, risk-taking, self-control, and reward.

\subsection*{Labour economics}

Another very important field is \textbf{labour economics}.

This is the study of jobs, wages, unemployment, education, skills, and the labour market.

Why do some jobs pay more than others?\\
What skills are in demand?\\
How does education affect earnings?\\
What happens when technology changes the workplace?

This area is especially relevant to students, because it connects directly to careers and future opportunities.

\subsection*{Environmental economics}

Then we have \textbf{environmental economics}.

This field studies pollution, climate change, energy, natural resources, and sustainability.

How should society reduce pollution?\\
What policies encourage cleaner energy?\\
How do we balance economic growth with environmental protection?

These are some of the biggest questions of our time.

\subsection*{Health economics, financial economics, development economics}

There are many more.

\textbf{Health economics} looks at healthcare systems, costs, and access.\\
\textbf{Financial economics} studies banking, markets, investment, and risk.\\
\textbf{Development economics} studies poverty, growth, education, and living standards around the world.

So economics is not one small box.\\
It is a wide field with many exciting paths.

\section*{Why economics matters now more than ever}

We live in a world of big changes.

Students today are growing up in a time shaped by inflation, housing pressures, rapid technology, artificial intelligence, climate concerns, and changing job markets.

Economics helps make sense of all of that.

Without economics, a lot of headlines feel confusing.

With economics, you begin to see the logic behind what is happening.

You understand why interest rates rise.\\
Why supply shortages affect prices.\\
Why labour markets change.\\
Why government policies matter.

Economics gives you a framework for understanding the world.

And that is incredibly valuable.

\section*{Challenges in economics}

Now, to be fair, economics is not magic.

It is powerful, but it has challenges.

Why?

Because people are complicated.

Humans do not always behave logically.\\
They are emotional.\\
They are social.\\
They make mistakes.\\
They change their minds.

Also, the world changes constantly.

Technology changes.\\
Markets change.\\
Policies change.\\
Unexpected events happen.

So economics often deals with probabilities, trade-offs, and imperfect information.

That is why economists sometimes disagree.

But that does not make economics weak.\\
It makes economics honest.

It means the world is complicated, and serious questions deserve serious thinking.

\section*{Why students should study economics}

Now let us talk about the practical part.

Why should a student study economics?

First, economics teaches students how to think.

It builds critical thinking, problem-solving, reasoning, and analytical skills.

It teaches students how to look at evidence, ask better questions, and make structured arguments.

Second, economics is practical.

It helps people understand everyday decisions about spending, saving, borrowing, careers, and policy.

Third, economics is flexible.

It opens doors to many careers:
government, banking, finance, consulting, policy analysis, research, business, education, data analysis, and more.

Even when the job title is not “economist,” the skills are still highly valuable.

Employers want people who can analyze information, think clearly, solve problems, and communicate well.

Economics helps develop exactly those abilities.

\section*{A direct message to students}

So let me say something directly to the students here.

You do not need to have your whole future figured out right now.

You do not need to know every answer.

What matters is curiosity.

If you like asking questions like:

Why is life getting more expensive?\\
Why do people behave the way they do?\\
Why do some jobs pay more than others?\\
Why do some policies work and others fail?

Then economics may be a very good fit for you.

Economics helps you understand the world --- and that is useful no matter what career you choose.

\section*{A direct message to parents}

And to the parents:

Economics is a field that offers both practical value and long-term flexibility.

It helps students develop strong thinking skills, strong communication skills, and strong analytical skills.

It prepares them not just for one job, but for many possible paths.

And in a world that is constantly changing, that kind of adaptability matters a lot.

Also, after taking economics, your children may finally understand why you keep saying,\\
“We cannot buy everything.”

That alone may justify the degree.

[Pause for laughter]

\section*{Closing}

Let me finish with this.

Economics is not just about money.\\
It is about choices.\\
It is about trade-offs.\\
It is about incentives.\\
It is about behaviour.\\
It is about understanding how the world works.

It explains everyday life.\\
It explains big issues.\\
It helps us think more clearly.\\
And it opens doors for the future.

Once you begin to see the world through an economic lens, you realize something important:

Economics is everywhere.

In prices.\\
In jobs.\\
In family decisions.\\
In public policy.\\
In housing.\\
In education.\\
In human behaviour.

And that is exactly why it is worth studying.

Thank you very much.

[Pause]

I hope after tonight, when you hear the word \textbf{economics}, you do not think “boring.”

I hope you think:

“That is actually about life.”

\end{document}